\documentclass{ximera}
%nieuw 2 september 2026
\title{Vraagstukken elektrisch veld: ladingen op een lijn}
\author{Daan Lipkens}

\begin{document}
\begin{abstract}
    Oefeningen op de superpositie van elektrische velden van puntladingen die zich op één rechte lijn bevinden.
\end{abstract}
\maketitle

%———————————————————————————————————————
% OEF LIJN 1: netto veld buiten twee tegengestelde ladingen
\begin{exercise}
    Twee puntladingen bevinden zich op de $x$-as: $Q_{1}=+4{,}0 \cdot 10^{-6} \, \mathrm{C}$ bij $x=0{,}00 \, \mathrm{m}$ en $Q_{2}=-4{,}0 \cdot 10^{-6} \, \mathrm{C}$ bij $x=0{,}60 \, \mathrm{m}$.

    \begin{center}
        \begin{tikzpicture}[scale=1]
            \draw[gray, thick] (-0.5,0) -- (4.5,0);
            \draw[charge+] (0,0) circle (\R) node[scale=0.5]{$+$};
            \draw[charge-] (3,0) circle (\R) node[scale=0.5]{$-$};
            \filldraw (4,0) circle (1.2pt);
            \node[below] at (0,-0.15) {$Q_1$};
            \node[below] at (3,-0.15) {$Q_2$};
            \node[below] at (4,-0.15) {$P$};
            \node[above] at (1.5,0.1) {$0{,}60\,\mathrm{m}$};
            \node[above] at (3.5,0.1) {$0{,}30\,\mathrm{m}$};
        \end{tikzpicture}
    \end{center}

    Bereken de grootte en richting van het totale elektrisch veld in punt \textit{P} op $x = 0{,}90 \, \mathrm{m}$.

    \[ E_{\text{net}} = \answer[tolerance=0.05,onlineshowanswerbutton]{3.6} \cdot 10^{5} \, \mathrm{N/C} \]

    \begin{hint}
        Bereken eerst de velden $E_1$ en $E_2$ afzonderlijk, elk met hun eigen afstand tot $P$.
    \end{hint}
    \begin{hint}
        Bepaal voor elk veld apart de zin (denk aan het teken van de bronlading), en tel daarna algebraïsch op: alle ladingen en $P$ liggen immers op dezelfde lijn.
    \end{hint}

    \begin{solution}\nl
        \underline{Gegevens:}
        \begin{align*}
            & Q_{1} = +4{,}0 \cdot 10^{-6} \, \mathrm{C} & & Q_{2} = -4{,}0 \cdot 10^{-6} \, \mathrm{C} \\
            & r_{1} = 0{,}90 \, \mathrm{m} & & r_{2} = 0{,}30 \, \mathrm{m}
        \end{align*}

        \underline{Gevraagd:} $E_{\text{net}}$ in $P$\\

        \underline{Oplossing:}\\
        \textbf{Stap 1: Grootte van de afzonderlijke velden}
        \begin{align*}
            E_{1} &= k_e \frac{|Q_{1}|}{r_{1}^{2}} = \frac{8{,}99 \cdot 10^{9} \, \mathrm{\frac{N \cdot m^2}{C^2}} \cdot 4{,}0 \cdot 10^{-6} \, \mathrm{C}}{(0{,}90 \, \mathrm{m})^{2}} = 4{,}4 \cdot 10^{4} \, \mathrm{N/C} \\
            E_{2} &= k_e \frac{|Q_{2}|}{r_{2}^{2}} = \frac{8{,}99 \cdot 10^{9} \, \mathrm{\frac{N \cdot m^2}{C^2}} \cdot 4{,}0 \cdot 10^{-6} \, \mathrm{C}}{(0{,}30 \, \mathrm{m})^{2}} = 4{,}0 \cdot 10^{5} \, \mathrm{N/C}
        \end{align*}

        \textbf{Stap 2: Zin van de velden bepalen}\\
        We kiezen de x-as naar rechts. Vectoren die naar rechts wijzen, zijn positief. Vectoren die naar links wijzen, zijn negatief.
        $Q_1$ is positief: $\vec{E}_1$ wijst in $P$ weg van $Q_1$, dus naar rechts ($+x$).\\
        $Q_2$ is negatief: $\vec{E}_2$ wijst in $P$ naar $Q_2$ toe, dus naar links ($-x$).

        \textbf{Stap 3: Algebraïsch optellen} (rechts is positief)
        \begin{align*}
            E_{\text{net}} &= E_1 - E_2 \\
                            &= 4{,}4 \cdot 10^{4} \, \mathrm{N/C} - 4{,}0 \cdot 10^{5} \, \mathrm{N/C} \\
                            &= -3{,}6 \cdot 10^{5} \, \mathrm{N/C}
        \end{align*}
        Het resulterende veld heeft een grootte van $3{,}6 \cdot 10^{5} \, \mathrm{N/C}$ en wijst naar \textbf{links} ($-x$-richting): dus naar de ladingen toe.
    \end{solution}
\end{exercise}
%———————————————————————————————————————

%———————————————————————————————————————
% OEF LIJN 2: nulpunt tussen 2 positieve ladingen
\begin{exercise}
    Twee positieve puntladingen bevinden zich op de $x$-as: $Q_{1}=+2{,}0 \cdot 10^{-6} \, \mathrm{C}$ bij $x=0{,}00 \, \mathrm{m}$ en $Q_{3}=+8{,}0 \cdot 10^{-6} \, \mathrm{C}$ bij $x=1{,}20 \, \mathrm{m}$.

    Op welke plaats tussen de twee ladingen is het totale elektrisch veld gelijk aan nul? Geef je antwoord als de afstand $x$ vanaf $Q_1$.

    \[ x = \answer[tolerance=0.02,onlineshowanswerbutton]{0.40} \, \mathrm{m} \]

    \begin{hint}
        $Q_2$ heb je niet nodig.
    \end{hint}

    \begin{hint}
        Noem de afstand tot $Q_1$ gelijk aan $x$, dan is de afstand tot $Q_3$ gelijk aan $(1{,}20 \, \mathrm{m} - x)$.
    \end{hint}
    \begin{hint}
        Op de gezochte plaats moeten de velden van $Q_1$ en $Q_3$ exact even groot zijn (en tegengesteld gericht). Stel $E_1 = E_3$.
    \end{hint}


    \begin{solution}\nl
        \underline{Gegevens:}
        \begin{align*}
            & Q_{1} = +2{,}0 \cdot 10^{-6} \, \mathrm{C} & & Q_{3} = +8{,}0 \cdot 10^{-6} \, \mathrm{C} & & d = 1{,}20 \, \mathrm{m}
        \end{align*}

        \underline{Gevraagd:} $x$, de afstand vanaf $Q_1$ waar $E_{\text{net}} = 0$\\

        \underline{Oplossing:}\\
        Omdat beide ladingen positief zijn, wijzen hun velden tussen de ladingen in tegengestelde richtingen: het is dus mogelijk dat ze elkaar opheffen. In het gezochte punt geldt $E_1 = E_3$.
        Met $E_1=k_e \frac{Q_1}{r_1^{2}}=k_e \frac{Q_1}{x^{2}}$ en $E_3=k_e \frac{Q_3}{r_3^{2}}=k_e \frac{Q_3}{(d-x)^{2}}$. Hieruit volgt:
        \begin{align*}
            k_e \frac{Q_1}{x^{2}} &= k_e \frac{Q_3}{(d-x)^{2}} \\
            \frac{Q_1}{x^{2}} &= \frac{Q_3}{(d-x)^{2}} \\
            \left(\frac{d-x}{x}\right)^{2} &= \frac{Q_3}{Q_1} = \frac{8{,}0 \cdot 10^{-6} \, \mathrm{C}}{2{,}0 \cdot 10^{-6} \, \mathrm{C}} = 4 \\
            \frac{d-x}{x} &= 2
        \end{align*}
        Hieruit volgt $d - x = 2x$, dus $d = 3x$, en
        \[
            x = \frac{d}{3} = \frac{1{,}20 \, \mathrm{m}}{3} = 0{,}40 \, \mathrm{m}
        \]
        Het gezocht punt ligt dus op $0{,}40 \, \mathrm{m}$ van $Q_1$ (en dus $0{,}80 \, \mathrm{m}$ van $Q_3$), zoals verwacht dichter bij de kleinste lading.
    \end{solution}
\end{exercise}
%———————————————————————————————————————

%———————————————————————————————————————
% OEF LIJN 3: onbekende lading vinden (multi-question)
\begin{exercise}
    Op de $x$-as bevindt zich een puntlading $Q_{1}=+8{,}0 \cdot 10^{-6} \, \mathrm{C}$ bij $x=0{,}00 \, \mathrm{m}$. Een tweede, onbekende puntlading $Q_{2}$ bevindt zich bij $x=0{,}60 \, \mathrm{m}$. In punt \textit{P} op $x=0{,}20 \, \mathrm{m}$ blijkt het totale elektrisch veld gelijk aan nul te zijn.

    \begin{question}
        Welk teken moet $Q_2$ hebben opdat het veld in $P$ nul kan zijn?

        \wordChoice{\choice[correct]{Positief}\choice{Negatief}}

        \begin{hint}
            Het veld van $Q_1$ in $P$ wijst naar rechts (weg van $Q_1$). Het veld van $Q_2$ in $P$ moet dus naar links wijzen. Punt $P$ ligt links van $Q_2$: in welke richting wijst het veld van een positieve lading, gezien vanaf een punt links ervan?
        \end{hint}

        \begin{solution}\nl
            $Q_1$ is positief, dus wijst $\vec{E}_1$ in $P$ weg van $Q_1$: naar rechts. Om dit te compenseren, moet $\vec{E}_2$ in $P$ naar links wijzen. Punt $P$ ligt links van $Q_2$; een veld dat daar naar links (dus weg van $Q_2$) wijst, hoort bij een \textbf{positieve} lading $Q_2$.
        \end{solution}
    \end{question}

    \begin{question}
        Bereken de grootte van $Q_2$.

        \[ Q_2 = \answer[tolerance=1,onlineshowanswerbutton]{32} \cdot 10^{-6} \, \mathrm{C} \]

        \begin{hint}
            In $P$ moet gelden $E_1 = E_2$. De afstand van $Q_1$ tot $P$ is $0{,}20 \, \mathrm{m}$; de afstand van $Q_2$ tot $P$ is $0{,}40 \, \mathrm{m}$.
        \end{hint}

        \begin{solution}\nl
            \underline{Gegevens:}
            \begin{align*}
                & Q_1 = +8{,}0 \cdot 10^{-6} \, \mathrm{C} & & r_1 = 0{,}20 \, \mathrm{m} & & r_2 = 0{,}40 \, \mathrm{m}
            \end{align*}

            \underline{Gevraagd:} $|Q_2|$\\

            \underline{Oplossing:}\\
            Opdat $E_{\text{net}}(P) = 0$, moet gelden $E_1 = E_2$:
            \begin{align*}
                k_e \frac{Q_1}{r_1^{2}} &= k_e \frac{Q_2}{r_2^{2}} \\
                Q_2 &= Q_1 \cdot \left(\frac{r_2}{r_1}\right)^{2} \\
                    &= 8{,}0 \cdot 10^{-6} \, \mathrm{C} \cdot \left(\frac{0{,}40 \, \mathrm{m}}{0{,}20 \, \mathrm{m}}\right)^{2} \\
                    &= 8{,}0 \cdot 10^{-6} \, \mathrm{C} \cdot 4 \\
                    &= 32 \cdot 10^{-6} \, \mathrm{C}
            \end{align*}
        \end{solution}
    \end{question}
\end{exercise}
%———————————————————————————————————————

%———————————————————————————————————————
% OEF LIJN 4: totaal veld, P tussen twee gelijknamige ladingen
\begin{exercise}
    Twee positieve puntladingen bevinden zich op de $x$-as: $Q_{1}=+3{,}0 \cdot 10^{-6} \, \mathrm{C}$ bij $x=0{,}00 \, \mathrm{m}$ en $Q_{2}=+7{,}0 \cdot 10^{-6} \, \mathrm{C}$ bij $x=0{,}80 \, \mathrm{m}$.

    \begin{center}
        \begin{tikzpicture}[scale=1]
            \draw[gray, thick] (-0.5,0) -- (4.5,0);
            \draw[charge+] (0,0) circle (\R) node[scale=0.5]{$+$};
            \draw[charge+] (4,0) circle (\R) node[scale=0.5]{$+$};
            \filldraw (1.5,0) circle (1.2pt);
            \node[below] at (0,-0.15) {$Q_1$};
            \node[below] at (4,-0.15) {$Q_2$};
            \node[below] at (1.5,-0.15) {$P$};
            \node[above] at (0.75,0.1) {$0{,}30\,\mathrm{m}$};
            \node[above] at (2.75,0.1) {$0{,}50\,\mathrm{m}$};
        \end{tikzpicture}
    \end{center}

    Bereken de grootte en richting van het totale elektrisch veld in punt \textit{P} op $x = 0{,}30 \, \mathrm{m}$.

    \[ E_{\text{net}} = \answer[tolerance=0.05,onlineshowanswerbutton]{4.8} \cdot 10^{4} \, \mathrm{N/C} \]

    \begin{hint}
        Beide ladingen zijn positief: in het gebied tussen hen wijzen $\vec{E}_1$ en $\vec{E}_2$ in tegengestelde richting.
    \end{hint}

    \begin{solution}\nl
        \underline{Gegevens:}
        \begin{align*}
            & Q_{1} = +3{,}0 \cdot 10^{-6} \, \mathrm{C} & & Q_{2} = +7{,}0 \cdot 10^{-6} \, \mathrm{C} \\
            & r_{1} = 0{,}30 \, \mathrm{m} & & r_{2} = 0{,}50 \, \mathrm{m}
        \end{align*}

        \underline{Gevraagd:} $E_{\text{net}}$ in $P$\\

        \underline{Oplossing:}\\
        \textbf{Stap 1: Grootte van de afzonderlijke velden}
        \begin{align*}
            E_{1} &= k_e \frac{|Q_{1}|}{r_{1}^{2}} = \frac{8{,}99 \cdot 10^{9} \, \mathrm{\frac{N \cdot m^2}{C^2}} \cdot 3{,}0 \cdot 10^{-6} \, \mathrm{C}}{(0{,}30 \, \mathrm{m})^{2}} = 3{,}0 \cdot 10^{5} \, \mathrm{N/C} \\
            E_{2} &= k_e \frac{|Q_{2}|}{r_{2}^{2}} = \frac{8{,}99 \cdot 10^{9} \, \mathrm{\frac{N \cdot m^2}{C^2}} \cdot 7{,}0 \cdot 10^{-6} \, \mathrm{C}}{(0{,}50 \, \mathrm{m})^{2}} = 2{,}5 \cdot 10^{5} \, \mathrm{N/C}
        \end{align*}

        \textbf{Stap 2: Zin van de velden bepalen}\\
        $Q_1$ is positief: $\vec{E}_1$ wijst in $P$ weg van $Q_1$, dus naar rechts ($+x$).\\
        $Q_2$ is positief: $\vec{E}_2$ wijst in $P$ weg van $Q_2$, dus naar links ($-x$).

        \textbf{Stap 3: Algebraïsch optellen} (rechts is positief)
        \begin{align*}
            E_{\text{net}} &= E_1 - E_2 = 3{,}0 \cdot 10^{5} \, \mathrm{N/C} - 2{,}5 \cdot 10^{5} \, \mathrm{N/C} = 4{,}8 \cdot 10^{4} \, \mathrm{N/C}
        \end{align*}
        Het resulterende veld wijst naar \textbf{rechts} ($+x$-richting), naar de sterkste lading $Q_2$ toe.
    \end{solution}
\end{exercise}
%———————————————————————————————————————

%———————————————————————————————————————
% OEF LIJN 5: nulpunt tussen 2 negatieve ladingen
\begin{exercise}
    Twee negatieve puntladingen bevinden zich op een lijn: $Q_{1}=-3{,}0 \cdot 10^{-6} \, \mathrm{C}$ en $Q_{4}=-12{,}0 \cdot 10^{-6} \, \mathrm{C}$. De ladingen liggen $90 \, \mathrm{cm}$ uit elkaar.

    Waar is het totale elektrisch veld gelijk aan nul? Geef je antwoord als de afstand $x$ vanaf $Q_1$. (Negatief als het links van $Q_1$ ligt en positief als het rechts van $Q_1$ ligt.)

    \[ x = \answer[tolerance=0.02,onlineshowanswerbutton]{0.30} \, \mathrm{m} \]

    \begin{hint}
        Noem de afstand tot $Q_1$ gelijk aan $x$, dan is de afstand tot $Q_4$ gelijk aan $(0{,}90 \, \mathrm{m} - x)$. Stel $E_1 = E_4$.
    \end{hint}

    \begin{solution}\nl
        \underline{Gegevens:}
        \begin{align*}
            & Q_{1} = -3{,}0 \cdot 10^{-6} \, \mathrm{C} & & Q_{4} = -12{,}0 \cdot 10^{-6} \, \mathrm{C} & & d = 0{,}90 \, \mathrm{m}
        \end{align*}

        \underline{Gevraagd:} $x$, de afstand vanaf $Q_1$ waar $E_{\text{net}} = 0$\\

        \underline{Oplossing:}\\
        Omdat beide ladingen negatief zijn, wijzen hun velden tussen de ladingen in tegengestelde richting (elk naar de eigen lading toe): het is dus mogelijk dat ze elkaar opheffen ergens tussen de twee ladingen. In het gezochte punt geldt $|E_1| = |E_4|$:
        \begin{align*}
            k_e \frac{|Q_1|}{x^{2}} &= k_e \frac{|Q_4|}{(d-x)^{2}} \\
            \left(\frac{d-x}{x}\right)^{2} &= \frac{|Q_4|}{|Q_1|} = \frac{12{,}0 \cdot 10^{-6} \, \mathrm{C}}{3{,}0 \cdot 10^{-6} \, \mathrm{C}} = 4 \\
            \frac{d-x}{x} &= 2
        \end{align*}
        Hieruit volgt $d - x = 2x$, dus $d = 3x$, en
        \[
            x = \frac{d}{3} = \frac{0{,}90 \, \mathrm{m}}{3} = 0{,}30 \, \mathrm{m}
        \]
        Het gezochte punt ligt dus op $0{,}30 \, \mathrm{m}$ van $Q_1$, zoals verwacht dichter bij de kleinste lading.
    \end{solution}
\end{exercise}
%———————————————————————————————————————

%———————————————————————————————————————
% OEF LIJN 6: totaal veld, P buiten, tegengestelde ladingen
\begin{exercise}
    Twee puntladingen bevinden zich op de $x$-as: $Q_{1}=+2{,}0 \cdot 10^{-6} \, \mathrm{C}$ bij $x=0{,}00 \, \mathrm{m}$ en $Q_{2}=-6{,}0 \cdot 10^{-6} \, \mathrm{C}$ bij $x=0{,}50 \, \mathrm{m}$.

    \begin{center}
        \begin{tikzpicture}[scale=1]
            \draw[gray, thick] (-2.5,0) -- (5.5,0);
            \filldraw (-2,0) circle (1.2pt);
            \draw[charge+] (0,0) circle (\R) node[scale=0.5]{$+$};
            \draw[charge-] (5,0) circle (\R) node[scale=0.5]{$-$};
            \node[below] at (-2,-0.15) {$P$};
            \node[below] at (0,-0.15) {$Q_1$};
            \node[below] at (5,-0.15) {$Q_2$};
            \node[above] at (-1,0.1) {$0{,}20\,\mathrm{m}$};
            \node[above] at (2.5,0.1) {$0{,}50\,\mathrm{m}$};
        \end{tikzpicture}
    \end{center}

    Bereken de grootte en richting van het totale elektrisch veld in punt \textit{P} op $x = -0{,}20 \, \mathrm{m}$.

    \[ E_{\text{net}} = \answer[tolerance=0.1,onlineshowanswerbutton]{3.4} \cdot 10^{5} \, \mathrm{N/C} \] en staat naar  \wordChoice{    
                \choice[correct]{links}  
                \choice{rechts}  
                \choice{boven}
                \choice{onder}  
                }  
    

    \begin{hint}
        $P$ ligt links van beide ladingen: $\vec{E}_1$ wijst dan weg van $Q_1$ (dus naar links), terwijl $\vec{E}_2$ naar de negatieve lading $Q_2$ toe wijst (dus naar rechts).
    \end{hint}

    \begin{solution}\nl
        \underline{Gegevens:}
        \begin{align*}
            & Q_{1} = +2{,}0 \cdot 10^{-6} \, \mathrm{C} & & Q_{2} = -6{,}0 \cdot 10^{-6} \, \mathrm{C} \\
            & r_{1} = 0{,}20 \, \mathrm{m} & & r_{2} = 0{,}70 \, \mathrm{m}
        \end{align*}

        \underline{Gevraagd:} $E_{\text{net}}$ in $P$\\

        \underline{Oplossing:}\\
        \textbf{Stap 1: Grootte van de afzonderlijke velden}
        \begin{align*}
            E_{1} &= k_e \frac{|Q_{1}|}{r_{1}^{2}} = \frac{8{,}99 \cdot 10^{9} \, \mathrm{\frac{N \cdot m^2}{C^2}} \cdot 2{,}0 \cdot 10^{-6} \, \mathrm{C}}{(0{,}20 \, \mathrm{m})^{2}} = 4{,}5 \cdot 10^{5} \, \mathrm{N/C} \\
            E_{2} &= k_e \frac{|Q_{2}|}{r_{2}^{2}} = \frac{8{,}99 \cdot 10^{9} \, \mathrm{\frac{N \cdot m^2}{C^2}} \cdot 6{,}0 \cdot 10^{-6} \, \mathrm{C}}{(0{,}70 \, \mathrm{m})^{2}} = 1{,}1 \cdot 10^{5} \, \mathrm{N/C}
        \end{align*}

        \textbf{Stap 2: Zin van de velden bepalen}\\
        $Q_1$ is positief: $\vec{E}_1$ wijst in $P$ weg van $Q_1$, dus naar links ($-x$).\\
        $Q_2$ is negatief: $\vec{E}_2$ wijst in $P$ naar $Q_2$ toe, dus naar rechts ($+x$).

        \textbf{Stap 3: Algebraïsch optellen} (rechts is positief)
        \begin{align*}
            E_{\text{net}} &= -E_1 + E_2 = -4{,}5 \cdot 10^{5} \, \mathrm{N/C} + 1{,}1 \cdot 10^{5} \, \mathrm{N/C} = -3{,}4 \cdot 10^{5} \, \mathrm{N/C}
        \end{align*}
        Het resulterende veld heeft een grootte van $3{,}4 \cdot 10^{5} \, \mathrm{N/C}$ en wijst naar \textbf{links} ($-x$-richting).
    \end{solution}
\end{exercise}
%———————————————————————————————————————

%———————————————————————————————————————
% OEF LIJN 7: totaal veld, meerdelig (richting + grootte), P ver buiten
\begin{exercise}
    Twee positieve puntladingen bevinden zich op de $x$-as: $Q_{1}=+1{,}0 \cdot 10^{-6} \, \mathrm{C}$ bij $x=0{,}00 \, \mathrm{m}$ en $Q_{2}=+9{,}0 \cdot 10^{-6} \, \mathrm{C}$ bij $x=1{,}00 \, \mathrm{m}$. Punt \textit{P} ligt op $x=1{,}50 \, \mathrm{m}$.

    \begin{question}
        Naar waar wijst het totale elektrisch veld in $P$?

        \wordChoice{\choice[correct]{Naar rechts ($+x$)}\choice{Naar links ($-x$)}}

        \begin{hint}
            $P$ ligt rechts van beide (positieve) ladingen. In welke richting wijst een veld van een positieve lading, gezien vanaf een punt rechts ervan?
        \end{hint}

        \begin{solution}\nl
            Beide ladingen zijn positief en $P$ ligt rechts van beide. De velden van $Q_1$ én $Q_2$ wijzen dus allebei weg van hun lading, naar rechts. Het totale veld wijst dus ook naar \textbf{rechts}.
        \end{solution}
    \end{question}

    \begin{question}
        Bereken de grootte van het totale elektrisch veld in $P$.

        \[ E_{\text{net}} = \answer[tolerance=0.1,onlineshowanswerbutton]{3.3} \cdot 10^{5} \, \mathrm{N/C} \]

        \begin{hint}
            De afstand van $Q_1$ tot $P$ is $1{,}50 \, \mathrm{m}$, die van $Q_2$ tot $P$ is $0{,}50 \, \mathrm{m}$. Omdat beide velden dezelfde zin hebben, mag je de groottes gewoon optellen.
        \end{hint}

        \begin{solution}\nl
            \underline{Gegevens:}
            \begin{align*}
                & Q_{1} = +1{,}0 \cdot 10^{-6} \, \mathrm{C} & & Q_{2} = +9{,}0 \cdot 10^{-6} \, \mathrm{C} \\
                & r_{1} = 1{,}50 \, \mathrm{m} & & r_{2} = 0{,}50 \, \mathrm{m}
            \end{align*}

            \underline{Gevraagd:} $E_{\text{net}}$ in $P$\\

            \underline{Oplossing:}\\
            \begin{align*}
                E_{1} &= k_e \frac{|Q_{1}|}{r_{1}^{2}} = \frac{8{,}99 \cdot 10^{9} \, \mathrm{\frac{N \cdot m^2}{C^2}} \cdot 1{,}0 \cdot 10^{-6} \, \mathrm{C}}{(1{,}50 \, \mathrm{m})^{2}} = 4{,}0 \cdot 10^{3} \, \mathrm{N/C} \\
                E_{2} &= k_e \frac{|Q_{2}|}{r_{2}^{2}} = \frac{8{,}99 \cdot 10^{9} \, \mathrm{\frac{N \cdot m^2}{C^2}} \cdot 9{,}0 \cdot 10^{-6} \, \mathrm{C}}{(0{,}50 \, \mathrm{m})^{2}} = 3{,}2 \cdot 10^{5} \, \mathrm{N/C}
            \end{align*}
            Beide velden wijzen naar rechts, dus:
            \begin{align*}
                E_{\text{net}} = E_1 + E_2 = 4{,}0 \cdot 10^{3} \, \mathrm{N/C} + 3{,}2 \cdot 10^{5} \, \mathrm{N/C} \simeq 3{,}3 \cdot 10^{5} \, \mathrm{N/C}
            \end{align*}
        \end{solution}
    \end{question}
\end{exercise}
%———————————————————————————————————————

%———————————————————————————————————————
% OEF LIJN 8: nulpunt BUITEN het segment, tegengestelde ladingen
\begin{exercise}
    Twee puntladingen bevinden zich op de $x$-as: $Q_{1}=+5{,}0 \cdot 10^{-6} \, \mathrm{C}$ bij $x=0{,}00 \, \mathrm{m}$ en $Q_{2}=-20{,}0 \cdot 10^{-6} \, \mathrm{C}$ bij $x=0{,}60 \, \mathrm{m}$.

    \begin{question}
        Aan welke kant van $Q_1$ ligt het punt waar het totale elektrisch veld nul is?

        \begin{multipleChoice}
            \choice{Tussen $Q_1$ en $Q_2$.}
            \choice[correct]{Links van $Q_1$.}
            \choice{Rechts van $Q_2$.}
        \end{multipleChoice}

        \begin{hint}
            Ga voor elk van de drie gebieden na of de velden van $Q_1$ en $Q_2$ elkaar daar kunnen tegenwerken. Bedenk ook dat het nulpunt, indien het buiten ligt, altijd aan de kant van de zwakste lading ligt.
        \end{hint}

        \begin{solution}\nl
            Tussen de ladingen wijzen beide velden in dezelfde zin (van $Q_1$ weg, naar $Q_2$ toe): daar kunnen ze elkaar dus nooit opheffen. Rechts van $Q_2$ is $Q_2$ altijd het dichtst en het sterkst, dus domineert dat veld daar volledig. Enkel links van $Q_1$, aan de kant van de zwakste lading, kan de kleinere lading $Q_1$ - dankzij de kleinere afstand - hetzelfde veld opwekken als de grotere, verder gelegen lading $Q_2$.
        \end{solution}
    \end{question}

    \begin{question}
        Bereken de afstand $x$ van dat punt tot $Q_1$.

        \[ x = \answer[tolerance=0.02,onlineshowanswerbutton]{0.60} \, \mathrm{m} \]

        \begin{hint}
            De afstand van het gezochte punt tot $Q_2$ is dan $(0{,}60 \, \mathrm{m} + x)$, want $Q_2$ ligt nu verder weg dan $Q_1$.
        \end{hint}

        \begin{solution}\nl
            \underline{Gegevens:}
            \begin{align*}
                & Q_{1} = +5{,}0 \cdot 10^{-6} \, \mathrm{C} & & Q_{2} = -20{,}0 \cdot 10^{-6} \, \mathrm{C} & & d = 0{,}60 \, \mathrm{m}
            \end{align*}

            \underline{Gevraagd:} $x$, de afstand links van $Q_1$ waar $E_{\text{net}} = 0$\\

            \underline{Oplossing:}\\
            Op het gezochte punt geldt $E_1 = E_2$, met afstand $x$ tot $Q_1$ en $(d+x)$ tot $Q_2$:
            \begin{align*}
                k_e \frac{|Q_1|}{x^{2}} &= k_e \frac{|Q_2|}{(d+x)^{2}} \\
                \left(\frac{d+x}{x}\right)^{2} &= \frac{|Q_2|}{|Q_1|} = \frac{20{,}0 \cdot 10^{-6} \, \mathrm{C}}{5{,}0 \cdot 10^{-6} \, \mathrm{C}} = 4 \\
                \frac{d+x}{x} &= 2
            \end{align*}
            Hieruit volgt $d + x = 2x$, dus $d = x$, en
            \[
                x = d = 0{,}60 \, \mathrm{m}
            \]
            Het gezochte punt ligt dus $0{,}60 \, \mathrm{m}$ links van $Q_1$.
        \end{solution}
    \end{question}
\end{exercise}
%———————————————————————————————————————

%———————————————————————————————————————
% OEF LIJN 9: nulpunt BUITEN, coördinaat gevraagd (hoogste moeilijkheidsgraad)
\begin{exercise}
    Twee puntladingen bevinden zich op de $x$-as: $Q_{1}=+4{,}0 \cdot 10^{-6} \, \mathrm{C}$ bij $x=0{,}00 \, \mathrm{m}$ en $Q_{2}=-9{,}0 \cdot 10^{-6} \, \mathrm{C}$ bij $x=1{,}00 \, \mathrm{m}$.

    \begin{question}
        Aan welke kant ligt het punt waar het totale elektrisch veld nul is?

        \begin{multipleChoice}
            \choice{Tussen $Q_1$ en $Q_2$.}
            \choice[correct]{Links van $Q_1$.}
            \choice{Rechts van $Q_2$.}
        \end{multipleChoice}

        \begin{hint}
            Vergelijk met een vorige oefening: bij tegengestelde ladingen ligt het nulpunt steeds buiten het segment, aan de kant van de zwakste lading.
        \end{hint}
    \end{question}

    \begin{question}
        Bepaal de $x$-coördinaat van dat punt.

        \[ x = \answer[tolerance=0.02,onlineshowanswerbutton]{-2.00} \, \mathrm{m} \]

        \begin{hint}
            Noem de (positieve) afstand tot $Q_1$ gelijk aan $d$. De afstand tot $Q_2$ is dan $(1{,}00 \, \mathrm{m} + d)$. Reken eerst $d$ uit, en zet dat resultaat pas op het einde om naar een $x$-coördinaat.
        \end{hint}

        \begin{solution}\nl
            \underline{Gegevens:}
            \begin{align*}
                & Q_{1} = +4{,}0 \cdot 10^{-6} \, \mathrm{C} & & Q_{2} = -9{,}0 \cdot 10^{-6} \, \mathrm{C} & & x_2 = 1{,}00 \, \mathrm{m}
            \end{align*}

            \underline{Gevraagd:} de $x$-coördinaat waar $E_{\text{net}} = 0$\\

            \underline{Oplossing:}\\
            Noem $d$ de afstand van het gezochte punt tot $Q_1$ (links van $Q_1$). De afstand tot $Q_2$ is dan $(1{,}00 \, \mathrm{m} + d)$. Op het gezochte punt geldt $E_1 = E_2$:
            \begin{align*}
                \left(\frac{1{,}00 \, \mathrm{m} + d}{d}\right)^{2} &= \frac{|Q_2|}{|Q_1|} = \frac{9{,}0 \cdot 10^{-6} \, \mathrm{C}}{4{,}0 \cdot 10^{-6} \, \mathrm{C}} = 2{,}25 \\
                \frac{1{,}00 \, \mathrm{m} + d}{d} &= 1{,}5
            \end{align*}
            Hieruit volgt $1{,}00 \, \mathrm{m} + d = 1{,}5d$, dus $1{,}00 \, \mathrm{m} = 0{,}5d$, en
            \[
                d = 2{,}00 \, \mathrm{m}
            \]
            Dit punt ligt $2{,}00 \, \mathrm{m}$ links van $Q_1$ (die op $x=0{,}00 \, \mathrm{m}$ ligt), dus is de gezochte coördinaat
            \[
                x = -2{,}00 \, \mathrm{m}
            \]
        \end{solution}
    \end{question}
\end{exercise}
%———————————————————————————————————————

\end{document}